\section{Introduction}
\label{sec:intro}

In 2021, the United States maternal mortality rate reached 32.9 deaths per 100,000
live births---the highest level since 1965---with Black mothers dying at nearly three
times the rate of White mothers~\citep{Toy2023WSJ,Harris2023}. Thrombotic embolisms
and hemorrhage together accounted for more than 22\% of these
deaths~\citep{Trost2022}, underscoring the central role of the coagulation system in
maternal outcomes. Blood coagulation is a complex enzymatic cascade in which a small
tissue factor (TF) stimulus triggers a series of serine protease activation reactions
that ultimately produce thrombin, the central effector of
hemostasis~\citep{Mann2011,Butenas2002}. Thrombin cleaves fibrinogen into fibrin,
activates platelets, and amplifies its own production through positive feedback via
Factors~V, VIII, and XI, while simultaneously initiating its own shutdown through the
thrombomodulin--protein~C anticoagulant pathway~\citep{Hockin2002}. Pregnancy
profoundly alters this hemostatic balance, shifting toward a hypercoagulable state
that prepares for delivery~\citep{Hellgren2003,Grouzi2022,Bremme2003}:
procoagulant factors including fibrinogen, Factor~VIII, and von Willebrand factor
increase substantially across gestation, while natural anticoagulants such as
antithrombin and protein~S decline. Conditions such as polycystic ovary syndrome
(PCOS)---a hormonal and metabolic disorder affecting 6--12\% of reproductive-age
women~\citep{Azziz2004} that has been associated with secondary hemostatic changes
including elevated Factor~VIII and impaired fibrinolysis---and preeclampsia (PE)---a hypertensive disorder of pregnancy
characterized by endothelial dysfunction, platelet activation, and altered thrombin
generation that complicates 2--8\% of pregnancies
worldwide~\citep{Abalos2013}---further perturb this balance in ways that remain incompletely
characterized. Mathematical models of the coagulation cascade, from the original
Hockin--Mann stoichiometric framework~\citep{Hockin2002} through extensions
incorporating the thrombomodulin--protein~C pathway~\citep{Bravo2012} to the
Brummel-Ziedins 2012 (BZ2012) model with detailed APC-mediated Factor~Va inactivation
kinetics~\citep{BrummelZiedins2012}, have provided mechanistic insight into
patient-specific thrombin generation~\citep{Luan2007,Luan2010}, but applying these
models to pregnancy cohorts requires longitudinal datasets capturing the full
complement of coagulation factors, inhibitors, and thrombin generation assay (TGA)
measurements across gestation---precisely the type of data that is rare and expensive
to collect.

The fundamental barrier is sample size. Longitudinal pregnancy coagulation studies
typically enroll tens of patients with complete follow-up across multiple trimesters,
yielding datasets where the number of features~$p$ (coagulation factors, TGA
parameters, viscoelastic measurements) exceeds the number of patients~$n$. Rare complications compound this problem: small longitudinal cohorts inevitably
contain only a handful of patients with any given condition, far too few for
condition-specific statistical analysis. This $n < p$ regime fundamentally limits conventional approaches:
covariance matrices are rank-deficient, cross-validation is unreliable, and overfitting
is nearly inevitable. Synthetic data generation offers a path forward by augmenting
small real datasets with statistically faithful synthetic records, enabling downstream
analyses---mechanistic modeling, clinical comparisons, hypothesis generation for rare
complications---that would otherwise be
infeasible~\citep{Gonzales2023,Murtaza2024}. However, existing methods face
fundamental limitations in the small-cohort longitudinal setting. Sampling from a
fitted multivariate normal (MVN) distribution is singular when $n < p$ and must be
regularized~\citep{Ledoit2004}, introducing bias that distorts the joint distribution;
MVN also assumes Gaussian marginals, linear correlations, and provides no mechanism
for conditional generation---it cannot amplify a specific subpopulation while
preserving its signatures. More expressive models such as GANs and
VAEs~\citep{Goodfellow2014,Kingma2014,Xu2019CTGAN}, including recent longitudinal
extensions~\citep{Murtaza2024}, require substantially larger training sets and are
prone to mode collapse at small $n$ (we confirm this empirically for CTGAN at
$n{=}23$ in this work).

Stochastic Attention (SA) provides an alternative framework rooted in modern Hopfield
network theory~\citep{Ramsauer2021}---a continuous generalization of the associative
memory model recognized by the 2024 Nobel Prize in Physics---in which stored patient
profiles serve as memory patterns in a continuous energy landscape and Langevin
dynamics generates novel samples that interpolate between them via the attention
mechanism. We recently introduced a
multiplicity-weighted extension of the Hopfield energy that adds per-pattern weights
to the log-sum-exp term, enabling continuous interpolation between unconditioned
generation and targeted subpopulation amplification at inference time without
retraining~\citep{Varner2024SA}. Crucially, SA operates on the data manifold directly
rather than fitting a parametric distribution, making it well-suited to the $n < p$
regime. In prior work, we generated synthetic patient populations using MVN
distributions fitted independently to each clinical visit~\citep{Varner2023JuliaCon};
while this captured per-visit statistics, it assumed Gaussian marginals and generated
each visit independently rather than as part of a coherent longitudinal trajectory.

In this work, we adapted multiplicity-weighted SA for longitudinal clinical data and
introduced a four-level validation framework to rigorously assess synthetic patient
quality. We formulated a pipeline including concatenated-visit patterns, PCA
dimensionality reduction, and a direction-magnitude decomposition that preserves
anisotropic variance structure, with the multiplicity-weighted energy providing a
unified mechanism for both unconditioned generation and targeted amplification of rare
subgroups. We validated synthetic patients through four progressively stringent
levels---marginal plausibility, cross-visit covariance preservation, conditional
subgroup amplification, and mechanistic consistency with the BZ2012 ODE
model~\citep{BrummelZiedins2012}---and demonstrated that SA-generated patients passed
all four levels when applied to a 72-feature, 3-visit pregnancy coagulation dataset
($K{=}23$ patients, including 3 PCOS and 5 PE), while MVN failed at levels~2
through~4, and that the BZ2012 model confirmed 83--92\% cloud overlap between real
and synthetic patients under physiological conditions.
